\chapter{Tolerances and defaults}\label{app:tolerances}

This appendix lists every tolerance-like parameter of the Python API with its default and
meaning. The values are taken from the signatures of the installed package and from the kernel
headers. Background on the tolerance model is in \cref{sec:concepts-tolerances}. All lengths are
in model units; ``model size'' means the size of the bounding box of the input.

\section{Kernel constants}

\begin{tabularx}{\textwidth}{@{}llX@{}}
\toprule
\textbf{C++ constant} & \textbf{Value} & \textbf{Meaning} \\
\midrule
\code{nf::tol::kLength} & $10^{-10}$ & Model-space lengths: zero lengths, coincident points,
  parallel lines. Also \code{nf::kTolerance} (alias). \\
\code{nf::tol::kParametric} & $10^{-10}$ & Parameter differences: equal knots, multiplicities,
  snapping. \\
\code{nf::tol::kAngular} & $10^{-9}$ & Angles in radians, dot products of unit vectors. \\
\bottomrule
\end{tabularx}

The class \apiref{Tolerance} bundles the three values for one operation:
\code{Tolerance(length=1e-10, parametric=1e-10, angular=1e-9)}. Functions that take a
\code{Tolerance} default to these values. Weights are compared relatively.

\section{Parameters by area}

In the following tables, \emph{Tol.} denotes a \code{Tolerance} object with the kernel defaults.
A default of 0 means that the value is derived automatically, as described in the last column.

\subsection{NURBS operations}
\begin{longtable}{@{}p{0.33\textwidth}p{0.14\textwidth}p{0.47\textwidth}@{}}
\toprule
\textbf{Function / parameter} & \textbf{Default} & \textbf{Meaning} \\
\midrule
\endhead
\code{NurbsCurve.remove_knot}, \code{try_remove_knot}: \code{tolerance} & $10^{-10}$ & Maximum
  Cartesian deviation allowed by knot removal. \\
\code{NurbsCurve.reduce_degree}: \code{tolerance} & $10^{-10}$ & Maximum deviation of the
  degree-reduced curve; \code{GeometryError} if not met. \\
\code{NurbsSurface.}\allowbreak\code{remove_knot_u/v}, \code{try_remove_knot_u/v}, \code{reduce_degree_u/v}:
  \code{tolerance} & $10^{-10}$ & As for curves, over all rows or columns. \\
\code{join}, \code{join_u}, \code{join_v}: \code{tolerance} & $10^{-10}$ & Maximum gap between
  the end of the first and the start of the second curve or surface. \\
\bottomrule
\end{longtable}

\subsection{Primitives, construction and modelling}
\begin{longtable}{@{}p{0.33\textwidth}p{0.14\textwidth}p{0.47\textwidth}@{}}
\toprule
\textbf{Function / parameter} & \textbf{Default} & \textbf{Meaning} \\
\midrule
\endhead
\code{make_helix}: \code{tolerance} & $10^{-7}$ & Deviation of the C2 quintic from the exact
  helix or spiral. \\
\code{SweepOptions}: \code{tolerance} & $10^{-6}$ & Deviation from the ideal swept surface,
  estimated at test points between the placed sections. \\
\code{SweepOptions}: \code{max_sections} & 1025 & Maximum number of placed sections;
  \code{GeometryError} if the tolerance is not met. \\
\code{coons}: \code{tolerance} & $10^{-7}$ & Corner agreement of the boundary curves and
  approximation of rational boundaries by non-rational ones. \\
\code{gordon}: \code{tolerance} & $10^{-7}$ & Consistency of the curve network (intersection
  parameters) and approximation of rational curves. \\
\code{offset_curve}: \code{tolerance} & $10^{-7}$ & Deviation of the fitted cubic offset. \\
\code{offset_surface}: \code{tolerance} & $10^{-7}$ & Deviation of the fitted quintic offset. \\
\code{fillet_curves}: \code{tolerance} & $10^{-7}$ & Fit tolerance of the auxiliary offset
  curves (start values only; the fillet arc is exact). \\
\code{fillet_corners}: \code{tolerance} & $10^{-7}$ & As for \code{fillet_curves}. \\
\code{fillet_surfaces}: \code{tolerance} & $10^{-6}$ & Accuracy of the fillet surface,
  contact curves and p-curves, verified between the sections. \\
\code{fillet_surfaces}: \code{offset_tolerance} & 0 & Fit tolerance of the auxiliary offset
  surfaces; 0 means \code{tolerance}/10. \\
\code{fillet_surfaces}: \code{max_sections} & 1025 & Maximum number of sections. \\
\bottomrule
\end{longtable}

\subsection{Fitting and optics}
\begin{longtable}{@{}p{0.33\textwidth}p{0.14\textwidth}p{0.47\textwidth}@{}}
\toprule
\textbf{Function / parameter} & \textbf{Default} & \textbf{Meaning} \\
\midrule
\endhead
\code{approximate_curve_tol}: \code{tolerance} & required & Maximum distance of the data
  points from the curve; the fewest control points meeting it are used. \\
\code{asphere_surface}, \code{asphere_profile}: \code{tolerance} & $10^{-7}$ & Sag error of
  the fitted profile, checked at 8 points per knot span. Pure conics are exact. \\
\code{zernike_surface}, \code{xy_polynomial_surface}: \code{tolerance} & $10^{-7}$ & Sag error,
  checked at $8\times8$ points per knot rectangle. \\
\code{zernike_surface}, \code{xy_polynomial_surface}: \code{half_width} & 0 & Half width of the
  square domain; 0 means \code{norm_radius}. \\
\code{sag_surface}: \code{tolerance} & $10^{-7}$ & Sag error of the graph surface. \\
\code{sag_surface}: \code{max_spans} & 256 & Maximum number of knot spans per direction. \\
\bottomrule
\end{longtable}

\subsection{Analysis and integration}
\begin{longtable}{@{}p{0.33\textwidth}p{0.14\textwidth}p{0.47\textwidth}@{}}
\toprule
\textbf{Function / parameter} & \textbf{Default} & \textbf{Meaning} \\
\midrule
\endhead
\code{closest_point}, \code{closest_point_local}: \code{tolerance} & Tol. & Convergence of the
  Newton iteration (length and parametric). \\
\code{intersect} (curve/curve, curve/surface, curve/plane): \code{tolerance} & Tol. & Accuracy
  of the intersection points. \\
\code{SurfaceIntersection}\allowbreak\code{Options}: \code{tolerance} & Tol. & Point accuracy of the marching. \\
\code{SurfaceIntersection}\allowbreak\code{Options}: \code{chord_tolerance} & 0 & Maximum deviation of the traced
  polyline from the true curve, controls the step; 0 means $10^{-4}$ of the model size. \\
\code{SurfaceIntersection}\allowbreak\code{Options}: \code{fit_tolerance} & 0 & Accuracy of the fitted 3D curve and
  p-curves, verified between the fit points; 0 means $10^{-7}$ of the model size. \\
\code{length}, \code{arc_length}, \code{parameter_at_length}: \code{rel_tol} & $10^{-10}$ &
  Relative accuracy of the adaptive integration. \\
\code{area} (surface): \code{rel_tol} & $10^{-10}$ & Relative accuracy. \\
\code{integrate}: \code{rel_tol}, \code{abs_tol} & $10^{-10}$, 0 & Relative and absolute accuracy
  of the Gauss--Kronrod integral. \\
\bottomrule
\end{longtable}

\subsection{Trimming, shells and Booleans}
\begin{longtable}{@{}p{0.33\textwidth}p{0.14\textwidth}p{0.47\textwidth}@{}}
\toprule
\textbf{Function / parameter} & \textbf{Default} & \textbf{Meaning} \\
\midrule
\endhead
\code{TrimmedSurface(...)}: \code{tolerance} & Tol. & Validation of the trim loops: they must be
  closed (within \code{length}, in parameter units) and lie in the parameter domain. \\
\code{TrimmedSurface.area}: \code{rel_tol} & $10^{-9}$ & Relative accuracy (Green's theorem). \\
\code{TrimmedSurface.}\allowbreak\code{boundary_curves}: \code{tolerance} & $10^{-7}$ & Fit tolerance of the 3D
  boundary curves where they cannot be exact. \\
\code{split}: \code{tolerance} & Tol. & Tolerances of the surface/surface intersection that
  produces the cut curves. \\
\code{Shell.sew}: \code{tolerance} & $10^{-6}$ & Two boundary curves form one edge if each lies on
  the other within this distance. \\
\code{Shell.area}, \code{volume}, \code{mass_properties}: \code{rel_tol} & $10^{-8}$ & Relative
  accuracy of the surface integrals. \\
\code{boolean}: \code{tolerance} & $10^{-6}$ & Contacts and overlaps thinner than this are
  treated as contacts (tangential contacts, coincident faces). \\
\code{point_in_shell}: \code{tolerance} & $10^{-6}$ & Geometric tolerance of the ray casting
  (rays that hit a face boundary or graze a face are discarded; the majority of several valid
  rays decides). \\
\bottomrule
\end{longtable}

\subsection{Sections and hatching}
\begin{longtable}{@{}p{0.33\textwidth}p{0.14\textwidth}p{0.47\textwidth}@{}}
\toprule
\textbf{Function / parameter} & \textbf{Default} & \textbf{Meaning} \\
\midrule
\endhead
\code{SectionOptions}: \code{join_tolerance} & $10^{-6}$ & Ends of neighbouring section pieces
  closer than this are joined (the sewing tolerance). \\
\code{SectionOptions}: \code{fit_tolerance} & 0 & Accuracy of traced section curves on general
  faces; 0 means $10^{-9}$ of the model size. Isoparametric sections are exact. \\
\code{critical_heights}, \code{is_extrusion_band}: \code{tolerance} & Tol. & Classification of
  heights. \\
\code{hatch}: \code{min_length} & 0 & Hatch segments shorter than this are dropped. \\
\code{hatch}: \code{tolerance} & 0 & Contacts of a loop with a hatch line without crossing do not
  split the line; 0 means $10^{-9}$ of the model size. \\
\code{hatch}: \code{contour_tolerance} & 0 & Fit tolerance of the contour curves; 0 means
  $10^{-7}$ of the model size. \\
\bottomrule
\end{longtable}

\subsection{STEP}
\begin{longtable}{@{}p{0.33\textwidth}p{0.14\textwidth}p{0.47\textwidth}@{}}
\toprule
\textbf{Function / parameter} & \textbf{Default} & \textbf{Meaning} \\
\midrule
\endhead
\code{write_step}, \code{to_step}: \code{uncertainty} & $10^{-6}$ & Written as \code{UNCERTAINTY_}\allowbreak\code{MEASURE_WITH_UNIT}; readers such as OpenCASCADE use it as the initial vertex
  and edge tolerance. Fitted edge curves are written within 1/100 of the sewing tolerance. \\
\code{write_step}, \code{to_step}: \code{unit} & \code{"mm"} & Declared length unit
  (\code{mm}, \code{cm}, \code{m}, \code{inch}, \code{um}). \\
\code{read_step}, \code{from_step} (and \code{_named}): \code{tolerance} & $10^{-6}$ & Sewing
  tolerance in target units; trim curves are fitted to a tenth of it. \\
\code{read_step}, \code{from_step}: \code{unit} & \code{"mm"} & Target unit of the conversion. \\
\bottomrule
\end{longtable}

\subsection{Meshing (\code{nurbforge.mesh})}
\begin{longtable}{@{}p{0.33\textwidth}p{0.14\textwidth}p{0.47\textwidth}@{}}
\toprule
\textbf{Function / parameter} & \textbf{Default} & \textbf{Meaning} \\
\midrule
\endhead
\code{GmshSession.mesh}: \code{size} & None & Global maximum element size; None keeps gmsh's
  default. \\
\code{GmshSession.mesh}: \code{order}, \code{dim} & 1, 2 & Element order and dimension. \\
\code{GmshSession.mesh}: \code{curvature} & None & Elements per $2\pi$ of curvature (at least
  1); None disables curvature-based sizes. \\
\code{GmshSession.mesh}: \code{heal} & True & Remove duplicate entities before meshing. \\
\code{GmshSession.refine}: \code{distance} & 0 & Distance over which the local size grows to the
  global size; 0 refines only the faces themselves. \\
\code{bem_mesh}: \code{size}, \code{order}, \code{curvature}, \code{refine} & None, 1, None, None
  & As for \code{GmshSession}; \code{refine} maps labels or surface tags to sizes. \\
\bottomrule
\end{longtable}

\subsection{LLM interface (\code{nurbforge.llm})}
\begin{longtable}{@{}p{0.33\textwidth}p{0.14\textwidth}p{0.47\textwidth}@{}}
\toprule
\textbf{Function / parameter} & \textbf{Default} & \textbf{Meaning} \\
\midrule
\endhead
\code{GeometryAgent}: \code{max_steps} & 20 & Maximum model turns (tool-call rounds) per
  request. \\
\code{GeometryAgent}: \code{verify_rounds} & 1 & Verification rounds after each answer. \\
\code{OllamaClient}: \code{timeout} & 300\,s & HTTP timeout per request. \\
\code{verify.check}: \code{tolerance} & None & None means $10^{-6}$ relative, at least $10^{-9}$
  absolute. \\
Measured properties (\code{inspect}) & $10^{-8}$ & Relative accuracy of lengths, areas and
  volumes. \\
\code{constraints.solve}: \code{tolerance}, \code{max_iterations} & $10^{-8}$, 40 &
  Levenberg--Marquardt stopping criteria on the relative target residuals. Results with residuals
  up to $10^{-6}$ count as converged. \\
\bottomrule
\end{longtable}
